% Actual analytic application, pages 31--38.
% Source exposition only; no new formal declarations are claimed.
\PaperPage{31}{4.1 The actual Riemann xi function}
\label{act:xi}
The application requires a fixed analytic function, including its values at
the apparent singularities of the usual completion. Write $\Lambda_0$ for
Mathlib's entire pole-subtracted completion, named
\texttt{completedRiemannZeta\textsubscript{0}} in the mathematical
description. The source definition is exactly
\begin{equation}\label{act:xi-normalization}
 \xi(s)=\frac{1+s(s-1)\Lambda_0(s)}{2},\qquad s\in\mathbb C.
\end{equation}
This formula defines an entire function without introducing removable
values by convention. At the two exceptional arguments it immediately gives
\begin{equation}
 \xi(0)=\xi(1)=\tfrac12.
\end{equation}
If $\Lambda$ denotes the meromorphic completed zeta function, the retained
comparison theorem identifies the familiar expression on its valid domain:
\begin{equation}
 \xi(s)=\frac{s(s-1)}2\Lambda(s)\qquad(s\ne0,1).
\end{equation}
Keeping these domains explicit prevents a totalized division operation in
the formal language from being mistaken for analytic evaluation at a pole.
It also fixes the normalization used later for the logarithmic derivative.

\begin{proposition}[Retained literal RH binding]\label{act:rh-binding}
For this particular $\xi$, Mathlib's proposition
\texttt{RiemannHypothesis} is equivalent to
\begin{equation}\label{act:xi-rh}
 \forall s\in\mathbb C,\qquad
 \xi(s)=0\ \Longrightarrow\ \operatorname{Re}s=\tfrac12.
\end{equation}
\end{proposition}
The source proves this equivalence by matching the zeros of $\xi$ with
the nontrivial zeros in Mathlib's definition. The latter excludes the
negative even integers and the exceptional argument $1$. The factors in
the completion are checked at the relevant arguments; the nonvanishing
theorems also rule out extraneous zeros. Thus the endpoint concerns the
ordinary Riemann hypothesis, rather than a proposition chosen to resemble
it or a model zero predicate.

This binding is a retained theorem from \texttt{ActualXi.lean}, and the
entire normalization comes from \texttt{AWPIRH/XiCore.lean}~\cite{source}.
We do not reprove the completion and nonvanishing infrastructure here.
Subsequent implications always return to this same proposition. In
particular, changing a moment convention or constructing an abstract
positive operator cannot silently change the analytic function to which
the conclusion applies. The role of the following pages is to exhibit
the exact connections that make such a return legitimate.
An overall nonzero constant would leave the logarithmic derivative
unchanged, but retaining the exact normalization makes the source
correspondence unambiguous throughout.

\PaperPage{32}{4.2 Entire descent through the centered square}
\label{act:descent}
The reflection equation of the actual completion places the center of
symmetry at one half. Introduce the centered function $F$ and the descended
function $\Theta$ by
\begin{equation}
 \xi(1-s)=\xi(s),\qquad F(w)=\xi(\tfrac12+w),\qquad F(-w)=F(w),
\end{equation}
\begin{equation}\label{act:theta-definition}
 \Theta(z)=F(\sqrt z),\qquad \Theta(w^2)=F(w).
\end{equation}
Here the square root in the definition is a chosen pointwise square root.
The second identity is independent of that choice: the two roots differ
only by sign, and $F$ is even. Neither this observation nor the notation
asserts that a global holomorphic square root exists.

\begin{proposition}[Retained entire descent]\label{act:theta-entire}
The function $\Theta$ is entire, satisfies
\eqref{act:theta-definition}, and has $\Theta(0)\ne0$.
\end{proposition}
For completeness, the analytic mechanism can be described locally. Near
a nonzero point one selects a holomorphic square root after normalizing
the argument into a neighborhood of $1$. Evenness identifies the resulting
composition with the pointwise definition. At zero, continuity and the
removable-singularity theorem complete the extension. Equivalently, an
even Taylor expansion descends by replacing $w^{2n}$ with $z^n$.
The source implements the local-root and removable-singularity route in
\texttt{ThetaEntire.lean}; origin nonvanishing is supplied separately.

The actual-function development also supplies conjugation symmetry:
\begin{equation}\label{act:theta-conjugation}
 \Theta(\overline z)=\overline{\Theta(z)},\qquad
 \Theta(z)=0\ \Longrightarrow\ \Theta(\overline z)=0.
\end{equation}
Its implications differ from those of evenness. Evenness identifies the
two centered xi arguments over a single square; conjugation pairs the
upper and lower half-plane zeros of the descended function. Neither
property alone places all zeros on a line.

The exact geometric translation is
\begin{equation}\label{act:theta-rh}
 \mathrm{RH}\quad\Longleftrightarrow\quad
 \bigl[\Theta(z)=0\ \Longrightarrow\ z\in(-\infty,0]\bigr].
\end{equation}
Indeed, a centered critical-line point has the form $it$, whose square
is $-t^2$. Conversely, a square is real and nonpositive precisely when
its square root has zero real part. Applying the centered-square identity
in both directions proves the translation, using the literal RH binding
already fixed. Origin nonvanishing subsequently removes the endpoint zero
from the actual divisor.
Away from that endpoint, the square map has nonzero derivative. Local
zero orders are therefore preserved when descending either member of
an opposite pair. Pairing those members avoids counting the same
descended zero twice.

\PaperPage{33}{4.3 Exact moments and their positive normalization}
\label{act:moments}
Origin nonvanishing gives a disk on which the actual logarithmic derivative
is holomorphic. Define the complex moments before taking real parts:
\begin{equation}\label{act:moment-definition}
 R(z)=\frac{\Theta'(z)}{\Theta(z)},\qquad
 a_n=\frac{(-1)^n}{n!}R^{(n)}(0),\qquad m_n=\operatorname{Re}a_n.
\end{equation}
The real coordinate in this definition is convenient for Bernstein rows,
but does not discard an unaccounted imaginary part. The retained reality
theorem proves $a_n=\overline{a_n}$ at every order. Schwarz reflection of
$\Theta$ passes to its derivative, quotient, and iterated derivatives;
evaluation at the real point zero therefore yields
\begin{equation}\label{act:moment-reality}
 a_n=m_n\in\mathbb R\qquad(n\ge0).
\end{equation}
This argument uses neither RH nor a preliminary assertion that the
canonical nodes are real. It is also independent of the later all-order
power-sum theorem. These distinctions are essential when the moments
are used to prove a zero-location criterion.

Write $A_i$ for the canonical complex scales of the paired actual
divisor, with multiplicity. The source supplies the absolutely convergent
first trace and its positivity:
\begin{equation}\label{act:kappa-trace}
 \kappa:=m_0=\sum_i\operatorname{Re}(A_i^{-1})>0.
\end{equation}
The proof uses the actual product, unconditional low-height zero exclusion,
and nonempty divisor support. Specifically, the scale geometry gives
$\operatorname{Re}(A_i^{-1})>0$ for every index. Summability permits comparison
with one strictly positive term, proving strict positivity of the total.
It does not assume that every scale is positive real; that stronger
conclusion would already encode the desired zero geometry.

The normalized moments and their derivative calibration are consequently
\begin{equation}\label{act:normalized-moments}
 u_n=\frac{m_n}{\kappa^{n+1}},\qquad u_0=1,\qquad
 R^{(n)}(0)=(-1)^n n!\kappa^{n+1}u_n.
\end{equation}
All divisions here are justified by the established strict positivity.
The sequence $u$ is now fixed by the actual analytic function, rather
than selected to satisfy a moment criterion. The accompanying modules
\texttt{ActualMomentReality.lean} and
\texttt{ActualCalibrationProgress.lean} supply the retained inputs~\cite{source}.
They discharge reality and normalization before either forward route
assumes $H_{00}$. The remaining question is whether this particular real
sequence has positive Bernstein rows at every order.
Reality itself supplies no such positivity: a real sequence can have
negative moments or negative finite differences. The alternating sign
convention fixes precisely which differences the later criterion tests.

\PaperPage{34}{4.4 Canonical nodes and exact power sums}
\label{act:calibration}
The paired divisor turns the moment sequence into spectral data without
presupposing a positive measure. With the complex scales $A_i$ from the
preceding page, define
\begin{equation}\label{act:nodes}
 q_i=(\kappa A_i)^{-1},\qquad
 \alpha_i=-A_i=-\frac1{\kappa q_i}.
\end{equation}
The canonical index retains repeated occurrences, so zero multiplicities
are preserved. Its calibration includes both inclusion and exhaustion:
\begin{equation}\label{act:zero-calibration}
 q_i\ne0,\quad \Theta(\alpha_i)=0,\qquad
 \Theta(z)=0\ \Longrightarrow\ \exists i\ (z=\alpha_i).
\end{equation}
Exhaustion is indispensable. A family consisting only of favorable zeros
could establish neither the spectral criterion nor the return to RH.
The source obtains it from the entire actual divisor and its opposite-pair
indexing, not from a finite list of numerically located zeros.

\begin{theorem}[Retained canonical power sums]\label{act:power-sums}
The actual nodes satisfy
\begin{equation}\label{act:power-sum-formula}
 \sum_i|q_i|<\infty,\qquad
 u_n=\sum_i q_i^{n+1}\quad(n\ge0),\qquad \sum_iq_i=1.
\end{equation}
The sums are convergent complex sums; their stated values are real.
\end{theorem}
The final mass identity is the case $n=0$. It does not imply positivity
of individual nodes, since cancellation remains possible in a complex
sum. Absolute summability, however, bounds each node and controls every
higher power, supplying the necessary convergence statements.

Here is the analytic content behind the all-order identification. Put
$b_i=A_i^{-1}$ and $M=1+\sum_i|b_i|$. On a common neighborhood of zero,
\begin{equation}
 |z|M<\tfrac12\quad\Longrightarrow\quad
 |1+zb_i|>\tfrac12,\qquad
 \left|\frac{b_i}{1+zb_i}\right|\le2|b_i|.
\end{equation}
The same summable family, multiplied by a constant depending on the
derivative order, majorizes every resolvent power. Locally uniform product
convergence identifies the logarithmic derivative with the resolvent sum;
holomorphic derivative interchange then produces all power sums.
Reality and the exact $\kappa$ normalization finish the calculation.

The module \texttt{CanonicalPowerSums.lean} supplies this construction;
\texttt{AnalyticClosure.lean} packages its fields as
\texttt{actualCanonicalCalibration}. Thus calibration is constructed
analytic data at the final endpoints. It remains an infinite mathematical
object, and its existence is neither a finite certificate nor a proof of
$H_{00}$.
Derivative interchange is justified separately at each finite order on
the same open neighborhood. Its majorant may depend on that order;
no uniform bound over all derivative orders is asserted or needed
for these identities.

\PaperPage{35}{4.5 Ellipse geometry and the spectral radius}
\label{act:radius}
For a complex node $q$, the quantity $|q|+|1-q|$ measures its position
relative to the two real foci zero and one. It is at least one by the
triangle inequality. Equality requires the two summands to have compatible
directions, giving the elementary characterization
\begin{equation}\label{act:segment-geometry}
 E(q):=|q|+|1-q|\ge1,\qquad
 E(q)=1\ \Longleftrightarrow\ q\in[0,1].
\end{equation}
The interval is regarded as a subset of $\mathbb C$. In particular,
real nodes outside the interval and genuinely nonreal nodes both give
strictly larger values. This geometry detects more than the imaginary
part alone.
For a fixed value greater than one, the level set is an ellipse with
those two foci. Thus the radius describes the largest ellipse reached
by the calibrated nodes. It is a radius for this affine Bernstein
geometry, not the usual supremum of the node moduli. Replacing it
by that simpler supremum would lose the contribution from the
factor involving one minus the node.

Define the radius of the actual calibrated family by
\begin{equation}\label{act:spectral-radius}
 \rho=\sup\bigl(\{1\}\cup\{E(q_i):i\in J\}\bigr).
\end{equation}
The inserted value one makes the lower bound explicit. Absolute
summability makes the family bounded, since
$E(q_i)\le1+2|q_i|$. Consequently the supremum is finite, and
\begin{equation}\label{act:radius-one}
 \rho=1\quad\Longleftrightarrow\quad \forall i\ (q_i\in[0,1]).
\end{equation}
For the actual calibration, this is exactly RH. If the nodes lie in the
interval, nonvanishing makes them positive; every calibrated zero
$-1/(\kappa q_i)$ is then negative real. Exhaustion and
\eqref{act:theta-rh} give RH. Conversely, RH places all descended zeros
on the negative axis and therefore makes each node positive real.
The identity $\sum_iq_i=1$ bounds each by one.

The connection with row variation is visible in the exact generating
polynomial. For $|x|\le1$, absolute summability justifies the formula
\begin{equation}\label{act:row-generating}
 P_N(x):=\sum_{j=0}^Nw_{N,j}x^j
       =\sum_iq_i\bigl(1-q_i+xq_i\bigr)^N,
 \qquad |P_N(x)|\le V_N.
\end{equation}
Taking absolute values before summing produces the useful bound
\begin{equation}\label{act:variation-upper}
 1\le V_N\le\sum_i|q_i|E(q_i)^N
                  \le\left(\sum_i|q_i|\right)\rho^N.
\end{equation}
This proves the upper exponential rate immediately. Establishing the
matching lower rate requires a separate argument, because complex
contributions to the generating polynomial can cancel. The next page
states precisely the retained theorem that resolves this issue.

\PaperPage{36}{4.6 The retained ordinary spectral limit}
\label{act:root-limit}
\begin{theorem}[Retained actual spectral closure]\label{act:spectral-limit}
For the actual normalized sequence and its canonical nodes,
\begin{equation}\label{act:ordinary-root-limit}
 \lim_{N\to\infty}V_N^{1/N}=\rho.
\end{equation}
This is an ordinary limit along positive integers.
\end{theorem}
The source constructs the required dominant extraction from the
calibration and then applies its spectral limit theorem. We retain that
result here; the following description identifies its obligations without
claiming to reproduce the full extraction proof. The case $\rho=1$ follows
directly from \eqref{act:variation-upper}. Suppose therefore that $\rho>1$.

Summability forces the nodes to approach zero outside finite sets, so
their ellipse radii approach one there. A maximizing node can therefore
be isolated in a finite region. Choosing a unit complex phase $x$ aligns
its two affine contributions. At this phase the row polynomial admits
the constructed decomposition
\begin{equation}\label{act:dominant-decomposition}
 P_N(x)=\sum_{\ell=1}^{d}c_\ell\lambda_\ell^N+T_N,
 \qquad |x|=1,\quad |\lambda_\ell|=\rho.
\end{equation}
The rates are distinct after grouping repeated nodes; at least one
coefficient is nonzero. Equal nodes have coefficients equal to their
positive integer multiplicity times the node, so grouping cannot erase
every dominant coefficient. The complement has a uniform strict bound:
\begin{equation}\label{act:tail-bound}
 |T_N|\le C r^N,\qquad 1\le r<\rho.
\end{equation}
This complement includes both the infinite tail and the finite nodes
whose affine rate is smaller than the maximum.

Distinct dominant phases can cancel at particular indices. The argument
therefore uses finite Vandermonde estimates over blocks of consecutive
indices, rather than asserting a pointwise lower bound for every
exponential sum. The strict tail estimate makes the remainder negligible.
The inequalities $|P_N(x)|\le V_N$ and monotonicity of $V_N$ transfer the
block lower estimates to the ordinary root limit. Those are the steps
for which the retained spectral infrastructure is needed.

Combining the limit with \eqref{act:radius-one} yields
\begin{equation}\label{act:subexp-rh}
 \bigl[\forall\varepsilon>0\ \exists A>0\ \forall N\ge0,
       \ V_N\le Ae^{\varepsilon N}\bigr]
       \quad\Longleftrightarrow\quad\mathrm{RH}.
\end{equation}
The exact source declarations are
\texttt{canonical\_actual\_root\_tendsto\_constructed} and
\texttt{canonical\_actual\_subexponential\_iff\_nativeRH} in
\texttt{AnalyticClosure.lean}~\cite{source}.
For the forward implication in this last criterion, each exponential
bound forces the limit to be at most its exponential base. Letting
that base decrease to one gives the required radius.

\PaperPage{37}{4.7 The first route and the shared converse}
\label{act:route-one}
The first route now combines finite cancellation algebra with the
retained actual spectral criterion. Recall that $H_{00}$ is semantic
admissibility of the existing cutoff-zero, coefficient-zero budget. Its
meaning includes every order; the finite notation for the candidate
does not restrict that quantifier. For the actual rows, the identities
proved earlier give
\begin{equation}\label{act:h00-chain}
 H_{00}\ \Longrightarrow\ C_n=0\ (n\ge0)
 \ \Longrightarrow\ V_{n+1}=V_n\ (n\ge0).
\end{equation}
The second implication follows from
$C_n=(n+1)(V_{n+1}-V_n)$ and the positivity of $n+1$.
Starting with the unconditional normalization $V_0=1$, induction yields
\begin{equation}
 H_{00}\quad\Longrightarrow\quad V_n=1\qquad(n\ge0).
\end{equation}
For each $\varepsilon>0$, choosing $A=1$ supplies the required
subexponential estimate:
\begin{equation}
 V_n=1\le e^{\varepsilon n}\qquad(n\ge0).
\end{equation}
Equation~\eqref{act:subexp-rh} therefore proves the first implication
$H_{00}\Rightarrow\mathrm{RH}$. Its mathematical content beyond finite
row algebra is exactly the constructed actual spectral closure. No
separate extraction or calibration hypothesis remains in the endpoint.
The proof also identifies the logical order: row positivity is available
directly from zero cancellation before the spectral implication is
invoked. The second route may therefore reuse that positivity without
assuming the first route's conclusion.

The reverse implication is shared by both routes. Assume RH. The actual
calibration and radius geometry place each $q_i$ in $[0,1]$. Substituting
the power sums into the finite Bernstein difference formula gives
\begin{equation}\label{act:positive-spectral-weights}
 w_{N,j}=\binom Nj\sum_iq_i^{j+1}(1-q_i)^{N-j}\ge0,
 \qquad 0\le j\le N.
\end{equation}
Absolute summability justifies exchanging the infinite node sum with
the finite difference sum. Each resulting term is nonnegative. The row
mass identity then gives $V_N=1$, hence zero cancellation and $H_{00}$.

We have consequently obtained the equivalence
\begin{equation}\label{act:first-equivalence}
 H_{00}\quad\Longleftrightarrow\quad\mathrm{RH}.
\end{equation}
The common reverse construction is exposed in the source as
\texttt{H00\_of\_nativeRH}. The equivalence does not discharge its
left-hand assertion: it characterizes RH by another universal statement.
In particular, checking any fixed collection of rows cannot supply the
induction hypothesis used above for all later degrees. A further theorem
would be needed to propagate finite data to every cancellation order.
This unresolved issue is separate from the validity of the two
conditional implications.

\PaperPage{38}{4.8 The second route and the actual RH endpoint}
\label{act:route-two}
The second forward construction applies the preceding measure--resolvent
argument to the actual moments. From $H_{00}$, finite row algebra gives
positive normalized Bernstein rows. The Hausdorff construction produces
a probability measure $\sigma$ on $[0,1]$, whose moments are exactly $u_n$.
Its transform is
\begin{equation}\label{act:actual-resolvent}
 S(z)=\kappa\int_{[0,1]}\frac{d\sigma(q)}{1+\kappa qz},
 \qquad \operatorname{Im}S(z)\le0\quad(\operatorname{Im}z>0).
\end{equation}
The actual derivative calibration \eqref{act:normalized-moments} identifies
the analytic germs of $S$ and $R$ at zero. Continuing the holomorphic
defect, rather than dividing at possible zeros, gives
\begin{equation}\label{act:actual-defect}
 \Theta'-S\Theta=0\quad\hbox{on the upper half-plane}.
\end{equation}
An upper-half-plane zero $a$ of multiplicity $m\ge1$ would have
$R(z)=m/(z-a)+B(z)$ with $B$ holomorphic nearby. At $z=a-i\varepsilon$
the principal part has imaginary part $m/\varepsilon>0$, contradicting
the sign after the bounded remainder is absorbed. Conjugate-zero symmetry
then makes every $\Theta$-zero real.

For an actual xi zero $s=\sigma_0+it$, centered-square descent therefore
implies
\begin{equation}\label{act:height-return}
 0=\operatorname{Im}(s-\tfrac12)^2
   =2(\sigma_0-\tfrac12)t,
 \qquad |t|>1.
\end{equation}
The last inequality is a retained unconditional zero-height theorem.
It forces $\sigma_0=1/2$, and Proposition~\ref{act:rh-binding} gives literal
RH. This closure needs no separate exclusion of positive-real
$\Theta$-zeros. Every analytic input is either constructed from $H_{00}$
or supplied for the fixed actual functions.
In particular, the measure is not posited independently of the input
sequence. Exact moment identification constructs the origin germ, and
origin nonvanishing suffices there. The continuation step is formulated
with an analytic defect precisely so that hypothetical zeros elsewhere
do not become unmentioned assumptions in the proof.

In namespace \texttt{RH\_HilbertLimit.TwoRoutes}, the four exact exported
declarations are
\begin{center}\small
\begin{tabular}{@{}l@{}}
\texttt{nativeRH\_of\_H00\_route\_one}\\
\texttt{H00\_iff\_nativeRH\_route\_one}\\
\texttt{nativeRH\_of\_H00\_route\_two}\\
\texttt{H00\_iff\_nativeRH\_route\_two}.
\end{tabular}
\end{center}
Their types are respectively $H00\to\texttt{RiemannHypothesis}$ and
$H00\leftrightarrow\texttt{RiemannHypothesis}$ for each route; the
formalization appendix records their verification context. The second
forward proof does not call the first RH endpoint, although both share
foundations and the reverse implication. The explanatory $L^2$ operator
model introduces no additional audited Lean result. Neither endpoint
proves $H_{00}$ unconditionally or constructs a terminating RH checker.
